% ============================================================
% question2.tex - 问题二的模型建立与求解
%
% 完整结构骨架与写法示范见 question1.tex，本文件仅提供占位。
% 后续 question3.tex / question4.tex 复制本文件并替换内容。
%
% 必须遵循的四小节顺序：
%   1. 方法对比与选择  (>= 2 种方法，含对比表 + 各自原理 + 选定理由)
%   2. 模型建立
%   3. 求解过程
%   4. 结果与分析      (>= 4 张图，每张 >= 100 字分析)
% 章节末尾必须加 \FloatBarrier。
% ============================================================

\section{问题二的模型建立与求解}

\subsection{方法对比与选择}

% 候选方法对比表（>= 2 种）+ 每种方法 100-200 字原理 + 选定理由
（内容）

\subsection{模型建立}

% 决策变量 + 目标函数 + 约束条件
（内容）

\subsection{求解过程}

% 算法步骤 + 参数设置 + 收敛条件
（内容）

\subsection{结果与分析}

% 结果表 + >= 4 张图（每图 >= 100 字分析）+ 灵敏度分析
（内容）

\FloatBarrier
